\section{Simulation Framework}
\label{sec:framework}

\subsection{Overview}

Fig.~\ref{fig:framework} shows the structure of the study. For each ED type and replication, the generator creates one patient list containing arrival times, triage levels, dispositions, pathways, all activity durations and the predicted attributes at every prediction-error level. Every policy is then run on the same list, so policy comparisons use common random numbers by construction. All model inputs are entered once in configuration files with their sources; the full specification and code are available (see the Data and Code Availability statement).

\begin{figure*}[!htbp]
\centering
\begin{adjustbox}{max width=\textwidth}
\begin{tikzpicture}[node distance=7mm, >=Stealth,
  blk/.style={draw, rounded corners=3pt, align=left, font=\footnotesize, fill=black!3, text width=4.9cm, minimum height=2.4cm, inner sep=4pt}]
\node[blk] (in) {\textbf{Inputs}\\[3pt]
  $\bullet$ ED types: acuity mix, load, configuration\\[1pt]
  $\bullet$ Published case mix and patient pathways\\[1pt]
  $\bullet$ Queue rules and routing rules};
\node[blk, right=of in] (sim) {\textbf{Simulation model}\\[3pt]
  $\bullet$ Patient generator (true and\\\phantom{$\bullet$} predicted attributes)\\[1pt]
  $\bullet$ Discrete-event model of the ED\\[1pt]
  $\bullet$ Common random numbers};
\node[blk, right=of sim] (an) {\textbf{Analyses}\\[3pt]
  $\bullet$ Policy fit and regret (RQ1)\\[1pt]
  $\bullet$ Value of prediction accuracy (RQ2)\\[1pt]
  $\bullet$ Scoring-rule fidelity and tuning (RQ3)\\[1pt]
  $\bullet$ Robustness checks};
\draw[->, thick] (in) -- (sim);
\draw[->, thick] (sim) -- (an);
\draw[->] (an.south) to[out=-110, in=-70] node[below, font=\scriptsize]{tuning of rule parameters} (sim.south);
\end{tikzpicture}
\end{adjustbox}
\caption{Structure of the study. The lower arrow indicates that the tuning of priority indices evaluates candidate rules by repeated simulation.}
\label{fig:framework}
\end{figure*}

\begin{figure*}[!htbp]
\centering
\begin{tikzpicture}[node distance=4mm and 5mm, >=Stealth,
  box/.style={draw, rounded corners=2pt, minimum height=7mm, align=center, font=\footnotesize, fill=black!3},
  res/.style={draw, rounded corners=2pt, minimum height=7mm, align=center, font=\footnotesize, fill=blue!6},
  q/.style={draw, dashed, rounded corners=2pt, align=center, font=\scriptsize}]
\node[box] (arr) {Arrival};
\node[res, right=of arr] (tr) {Triage\\(nurse)};
\node[box, right=of tr] (rg) {Registration};
\node[q, right=of rg] (rt) {Routing rule\\Main / second area};
\node[res, right=of rt] (sp) {Treatment space\\(bed or chair)};
\node[res, right=of sp] (as) {Physician\\initial assessment};
\node[res, below=9mm of as] (dx) {Pathway steps\\LB, XR, CT, DR};
\node[res, left=of dx] (ra) {Physician\\reassessment};
\node[box, left=of ra] (dec) {Decision};
\node[box, left=11mm of dec, yshift=4mm] (dd) {Discharge\\process};
\node[box, left=11mm of dec, yshift=-6mm] (bd) {Boarding\\(admitted)};
\node[box, left=of dd, yshift=-5mm] (dep) {Departure};
\draw[->] (arr) -- (tr); \draw[->] (tr) -- (rg); \draw[->] (rg) -- (rt); \draw[->] (rt) -- (sp); \draw[->] (sp) -- (as);
\draw[->] (as) -- node[right, font=\scriptsize]{if steps} (dx); \draw[->] (dx) -- (ra); \draw[->] (ra) -- (dec);
\draw[->] (as.south west) to[out=-150, in=30] node[below right, font=\scriptsize, pos=0.25]{no steps} (dec.north east);
\draw[->] (dec) -- (dd); \draw[->] (dec) -- (bd); \draw[->] (dd) -- (dep); \draw[->] (bd) -- (dep);
\node[font=\scriptsize, align=left, below=2mm of bd.south east, anchor=north west, xshift=-2mm] {Every resource (blue) has its own queue, ordered by the policy's queue rule;\\CTAS~1 patients are always served first.};
\end{tikzpicture}
\caption{Patient journey in the simulation model. Pathway steps follow the patient's sampled pathway in order (LB laboratory, XR X-ray, CT computed tomography, DR consultation).}
\label{fig:process}
\end{figure*}

\subsection{Process Model}

Figure~\ref{fig:process} shows the patient journey. A patient arrives, is triaged by a nurse (first-come-first-served, with CTAS~1 first), registers, and is routed to a treatment area. The patient then waits for a treatment space and, once in a space, for a physician initial assessment. Patients whose pathway includes diagnostic or consultation steps then go through those steps in the order of their pathway; each step uses a multi-server resource with its own queue (laboratory, X-ray, CT, consultant). After the last step a physician reassesses the patient and decides on disposition. Discharged patients go through a short discharge process; admitted patients board in the ED until an inpatient bed is available. Physicians work three 8-hour shifts; capacity changes at shift boundaries are non-preemptive.

The model has three configurations. In the \emph{pooled} configuration all patients share one area and one physician pool. In the \emph{fast-track} configuration a second area has its own spaces and physicians; a patient waiting for a fast-track space may take an idle Main space when no Main patient is waiting. In the \emph{split-flow} configuration the second area is a vertical area with chairs that are held only during physician encounters; patients wait for test results without occupying a chair, physicians are shared between the areas, and an admitted vertical patient boards in a Main bed. Main-area beds, and fast-track spaces, are held from the moment a patient enters them until departure.

\subsection{ED Types}

ED types are defined by three factors in a full factorial design. The \emph{acuity mix} has three levels (Table~\ref{tab:parameters}): a community mix equal to the pooled CTAS distribution of the 16 highest-volume EDs in Alberta, and tertiary and urgent-care-like mixes placed inside the between-site ranges reported for the same EDs, near their high- and low-acuity ends \cite{Holroyd2016_EDCaseMix}. These mixes are consistent with full-length studies of Canadian EDs: in three urban adult EDs in Calgary, one of them a regional cardiac, stroke and trauma centre, 27\%, 49\%, 19\% and 5\% of the visits triaged to CTAS levels 2--5 were at levels 2, 3, 4 and 5, respectively \cite{McRae2022_CrowdingRevisits}, a distribution between our community and tertiary mixes. Admission probabilities by CTAS level are the midpoints of the reported between-site ranges. The \emph{load} has three levels: the offered load of the busiest physician pool is \LoadLow, \LoadMedium{} or \LoadHigh. The \emph{configuration} is pooled, fast track or split flow.

\begin{table}[!htbp]
\centering
\caption{Acuity mixes, admission probabilities, CTAS targets and mean physician assessment times by CTAS level.}
\label{tab:parameters}
\footnotesize
\setlength{\tabcolsep}{3.5pt}
\begin{adjustbox}{max width=\linewidth}\begin{tabular}{lrrrrr}
\toprule
 & CTAS 1 & CTAS 2 & CTAS 3 & CTAS 4 & CTAS 5 \\
\midrule
Tertiary mix (\%) & 2.5 & 32.0 & 45.0 & 17.0 & 3.5 \\
Community mix (\%) & 1.0 & 18.0 & 44.0 & 30.0 & 7.0 \\
Urgent-care-like mix (\%) & 0.2 & 4.0 & 34.8 & 49.0 & 12.0 \\
Admission probability (\%) & 59.6 & 30.5 & 15.0 & 5.1 & 4.6 \\
Target, door to physician (min) & 15 & 15 & 30 & 60 & 120 \\
Mean initial assessment (min) & 40 & 30 & 25 & 18 & 12 \\
\bottomrule
\end{tabular}\end{adjustbox}
\end{table}

Capacities are derived from these inputs by a fixed rule, sized separately for the CTAS~1--3 and CTAS~4--5 streams, so that all three configurations of one acuity mix have the same total numbers of physicians per shift, treatment spaces and diagnostic servers. The base volume is \BaseVisits{} visits per day with a fixed diurnal arrival profile. Physicians are sized for an offered load of \PhysTargetMedium{} in each shift and stream; treatment spaces for an offered load of \SpaceSizing{} on the service-only holding time; diagnostic resources for \DiagTarget; and triage nurses for \TriageTarget{} at the busiest hour of the high load level. The arrival rate is then scaled so that the busiest physician pool has the target offered load. Because a dedicated fast track cannot borrow physicians from the Main area, a fast-track ED reaches a given load level with fewer visits than a pooled ED with the same staff. Table~\aref{tab:capacity} lists the resulting capacities and volumes, which range from \VisitsMin{} to \VisitsMax{} visits per day. Two further factors are varied one at a time on the medium-load types as a sensitivity analysis: boarding (mean decision-to-departure time of \BoardLowMean{} minutes in the core design and \BoardHighMean{} minutes in the high level) and diagnostic intensity (light, typical, heavy).

\subsection{Patient Generator}

Arrivals follow a non-homogeneous Poisson process with an hourly profile. Each patient receives a CTAS level from the acuity mix and an admission indicator from the CTAS-specific probability. The patient then draws a whole pathway rather than step-by-step transitions, which keeps the sequences realistic and avoids loops that do not occur in practice. Pathways come from a process-mining study of ED patient journeys that reported the \NPathways{} most frequent pathways, covering most visits \cite{Rismanchian2023_PatientJourney}. A pathway is drawn from those that end in the patient's disposition, with probability proportional to its published frequency multiplied by $\exp(\theta_l k)$, where $k$ is the number of diagnostic and consultation steps and $\theta_l$ a CTAS-specific tilt (from \TiltOne{} for CTAS~1 to \TiltFive{} for CTAS~5, and $0$ for CTAS~3). The diagnostic-intensity factor shifts every $\theta_l$ by $\pm$\DiagShiftHeavy. Activity durations are lognormal (Appendix~\aref{app:model}).

Each patient carries a six-attribute code (Appendix~\aref{app:coding}): A, CTAS level (1--5); B, service class (1 very long to 5 very short); C, admission (1 admitted, 2 discharged); D, diagnostics (1 yes, 2 no); E, resource use (1 three or more steps to 4 none); F, consultation (1 yes, 2 no). Attribute A is observed at triage. Attributes B--F are predicted. For each attribute with reference confusion matrix $M_{\text{ref}}$, the predicted class is drawn from row $p$ of
\begin{equation}
M(\alpha) = I + \alpha\,(M_{\text{ref}} - I),
\label{eq:alpha}
\end{equation}
where $p$ is the true class. Thus $\alpha=0$ gives perfect information and $\alpha=1$ the reference accuracy, which is \AccAdm\% for admission, \AccSvc\% for the five-class service time (errors fall mostly on neighbouring classes), \AccDiag\% for diagnostics and \AccCons\% for consultation. These accuracies are assumptions in the range reported for triage-time prediction models \cite{Hong2018_disposition,SUN2011_disposition}. Each patient's prediction at every $\alpha$ is obtained from the same uniform random number, so the comparison across $\alpha$ also uses common random numbers.

\subsection{Streaming Policies}

A policy combines a queue rule and, in the two-area configurations, a routing rule. The queue rule orders every queue in the ED: treatment spaces, physicians and diagnostic resources. Seven queue rules are compared:
\begin{itemize}
\item \textbf{FCFS}: first-come-first-served.
\item \textbf{CTAS}: strict priority by CTAS level, first-come-first-served within a level.
\item \textbf{APQ}: accumulating priority, equal to the elapsed wait divided by the CTAS target time; the patient with the largest value is served \cite{StanfordTaylorZiedins2014_APQ}.
\item \textbf{CTAS-SPT}: CTAS priority, then shortest predicted service (highest service class) within a level.
\item \textbf{CTAS-PD}: CTAS priority, then predicted discharge first within a level.
\item \textbf{SIM}: a similarity index (defined below) with target code $1$-$5$-$2$-$2$-$4$-$2$ and weights (\SimWA, \SimWB, \SimWC, \SimWD, \SimWE, \SimWF), which expresses ``acuity first, then quick and simple patients''; these values are a judgement-based setting of the kind used in practice \cite{GT-in-ED_HCMS2016}, and Section~\ref{sec:optimization} replaces them by optimized values.
\item \textbf{SIM-AGE}: SIM plus an ageing term of \SimAging{} per elapsed CTAS target time, which lets long-waiting patients move up.
\end{itemize}
Three routing rules decide which patients go to the second area: \textbf{CTAS streaming} (C) sends CTAS~4--5 patients; \textbf{prediction streaming} (P) sends CTAS~3--5 patients with predicted discharge and predicted service class 4 or 5; \textbf{similarity routing} (S) sends a CTAS~3--5 patient to the area whose target code is more similar to the patient's code (Main target $1$-$1$-$1$-$1$-$1$-$1$, second-area target $5$-$5$-$2$-$2$-$4$-$2$; ties to Main). A policy is written as queue rule/routing rule, for example CTAS-SPT/C. The pooled configuration has \NQueueRules{} policies and the two-area configurations \NPoliciesTwoArea. Two rules apply to all policies. CTAS~1 patients are served first at every queue. Ties are broken explicitly: first by the time the patient joined the queue, then by patient identifier; similarity scores are rounded to $10^{-9}$ before comparison so that equal scores tie exactly. We refer to CTAS/C as the \emph{default} policy because it reflects common practice.

\subsection{The Similarity-Based Priority Index and Its Fidelity}
\label{sec:propositions}

A priority index assigns each waiting patient a number computed from the patient's own attributes and serves the patient with the largest number. The similarity-based priority index used here needs only the coded triage attributes, a weighted sum and a sort, and its weights can be read directly, so it can be configured in an ED information system without optimization at run time. Let $p=(p_1,\dots,p_K)$ be a patient code with $p_k\in\{1,\dots,n_k\}$, $R_k=n_k-1$ the range of attribute $k$, $q$ a target code and $w\geq 0$ a weight vector. The weighted similarity of Offodile \cite{OFFODILE1990} is
\begin{equation}
S(p) = \sum_{k=1}^{K} w_k\left(1-\frac{|p_k-q_k|}{R_k}\right).
\label{eq:sim}
\end{equation}
The similarity index serves the patient with the largest $S$, with the safety and tie rules above. The following results describe when it reproduces a classical rule.

\begin{proposition}[Linear index]
\label{prop:linear}
If every attribute with $w_k>0$ has its target at an end of its range ($q_k\in\{1,n_k\}$), then $S$ is affine in $p$ and the rule is a static linear priority index that is monotone in every weighted attribute. If some weighted attribute has an interior target, $S$ is not monotone in that attribute, and the rule cannot coincide with any rule that is strictly monotone in it.
\end{proposition}


\begin{proposition}[Lexicographic dominance]
\label{prop:lex}
Let all weighted targets lie at an end of their ranges and order the attributes $k=1,\dots,K$. If
\begin{equation}
\frac{w_k}{R_k} > \sum_{j>k} w_j \quad \text{for every } k \text{ with } w_k>0,
\label{eq:lexcond}
\end{equation}
then the similarity index orders patients lexicographically by $(|p_1-q_1|,\dots,|p_K-q_K|)$, with first-come-first-served ties. If (\ref{eq:lexcond}) fails for some $k$ and all code combinations can occur, there exist two patients whom the similarity index does not order lexicographically.
\end{proposition}


Proposition~\ref{prop:lex} gives exact similarity-index representations of three classical rules. Weights $(1,0,\dots,0)$ give strict CTAS priority. Weights with $w_A/4>w_B$ and target $q_B=5$ give CTAS-SPT; we use $w_A=\RepSptWA$, $w_B=\RepSptWB$. Weights with $w_A/4>w_C$ and $q_C=2$ give CTAS-PD. The SIM weights do not satisfy (\ref{eq:lexcond}): $w_A/R_A=\SimWA/4$ is smaller than the sum of the other weights. Although SIM is often read as ``acuity first, then quick patients'', Proposition~\ref{prop:lex} shows that it is not CTAS-SPT; the experiments in Section~\ref{sec:results} measure how far it is from it.

\begin{proposition}[Two-target routing]
\label{prop:route}
Let a patient be routed to the second area if $S_{\text{second}}(p)>S_{\text{main}}(p)$, where the two scores use the same weights and targets that lie at opposite ends of every weighted attribute. Then
\begin{equation}
S_{\text{second}}(p)-S_{\text{main}}(p)=\sum_{k}\sigma_k\frac{w_k}{R_k}\,(2p_k-n_k-1),
\end{equation}
with $\sigma_k=\pm 1$, so routing is a linear threshold rule whose boundary passes through the centre of the code ranges. Weights on CTAS only reproduce CTAS streaming exactly, with CTAS~3 going to Main. Prediction streaming cannot be reproduced by any such weights.
\end{proposition}


Proofs are given in Appendix~\aref{app:proofs}, which adds two results. Proposition~4 shows that prediction streaming becomes representable when the routing rule compares $S_{\text{second}}-S_{\text{main}}$ with an offset $b$ instead of zero: with weights on service class and admission only, the rule equals prediction streaming if and only if $w_B<4w_C$ and $\max(w_C,w_B-w_C)\le b<w_C+w_B/2$. Proposition~5 notes that with all weights zero every score ties and the rule is FCFS, and with all weights zero and an ageing term SIM-AGE is APQ. The similarity indices with ageing and offset therefore contain FCFS, CTAS, CTAS-SPT, CTAS-PD, APQ and all three routing rules as special cases, but only under the stated conditions. We check the propositions by exhaustive enumeration of codes in the test suite and measure fidelity in the experiments.
